\chapter{Sicurezza Fisica e delle Infrastrutture}

\section{Panoramica}
La sicurezza fisica per i sistemi informativi ha due scopi principali:
\begin{enumerate}
    \item Prevenire danni all'infrastruttura fisica (hardware, edifici, energia elettrica).
    \item Prevenire l'accesso e l'uso improprio di tale infrastruttura, al fine di proteggere le informazioni in essa contenute.
\end{enumerate}

\section{Minacce Principali}

Le minacce fisiche possono essere raggruppate in tre macro-categorie generali:

\begin{itemize}
    \item \textbf{Ambientali e Naturali:} Includono eventi catastrofici estremi (terremoti, inondazioni) e condizioni ambientali locali inadeguate, come temperature fuori norma, umidità sbilanciata, incendi, danni da acqua, polvere e infestazioni.
    \item \textbf{Tecniche:} Relativamente all'infrastruttura di supporto. Comprendono anomalie dell'alimentazione elettrica (sottotensioni, blackout, picchi di sovratensione) e le interferenze elettromagnetiche (EMI) che disturbano l'elettronica.
    \item \textbf{Causate dall'Uomo:} Spesso le più imprevedibili, includono l'accesso fisico non autorizzato alle strutture, il furto di hardware e dati, il vandalismo e il sabotaggio intenzionale.
\end{itemize}

\section{Misure di Mitigazione e Prevenzione}

\begin{defbox}[Controlli e Prevenzione]
\begin{itemize}
    \item \textbf{Controlli Ambientali:} Sistemi HVAC per la regolazione del clima, rilevatori di fumo, sensori di allagamento e sistemi automatici per l'estinzione degli incendi.
    \item \textbf{Sicurezza Elettrica:} Gruppi di continuità (UPS) per brevi interruzioni e filtri di rete, affiancati da generatori a combustibile per sopperire a blackout prolungati.
    \item \textbf{Controllo degli Accessi:} Recinzioni, portinerie, serrature biometriche, badge tracciati, allarmi e videosorveglianza continua per impedire intrusioni fisiche.
\end{itemize}
\end{defbox}